%! Characteristics of Radical Functions — Unit 5, Lesson 5.0 — Exit Ticket
\documentclass[10pt]{article}
\usepackage{algebra2-article}
\usepackage{algebra2-boxes}
\newcommand{\TallMath}[1]{$\displaystyle #1\rule[-1.4em]{0pt}{3.2em}$}
\newcommand{\lgrid}[4]{%
  \draw[step=1, linegray!40, very thin] (#1,#3) grid (#2,#4);
  \draw[->, semithick, charcoal] (#1,0)--(#2,0) node[below right, font=\scriptsize]{$x$};
  \draw[->, semithick, charcoal] (0,#3)--(0,#4) node[left, font=\scriptsize]{$y$};}

\begin{document}

\pageheader{Unit 5, Lesson 5.0}{Exit Ticket}
\namedateperiod

\vspace{0.08in}

No notes. The graph below is $f(x) = \sqrt{x+1}-2$. Use it for questions 1 and 2.
\begin{center}
\begin{tikzpicture}[scale=0.44]
  \lgrid{-3}{10}{-3}{2}
  \begin{scope}
    \clip (-3,-3) rectangle (10,2);
    \draw[very thick, forest, domain=-1:10, samples=100, smooth] plot (\x,{sqrt((\x)+1)-2});
  \end{scope}
  \fill[forest] (-1,-2) circle (4pt) node[below left, font=\scriptsize]{$(-1,-2)$};
  \fill[charcoal] (0,-1) circle (3pt) node[left, font=\scriptsize]{$(0,-1)$};
  \fill[charcoal] (3,0) circle (3pt) node[above left, font=\scriptsize]{$(3,0)$};
  \fill[charcoal] (8,1) circle (3pt) node[below right, font=\scriptsize]{$(8,1)$};
\end{tikzpicture}
\end{center}

\begin{enumerate}[leftmargin=1.8em, itemsep=8pt]
  \item \textbf{Read the features.} Write the domain and range in interval notation.

        Domain: \blank{3.2cm} \quad Range: \blank{3.2cm} \quad Endpoint: \blank{2.0cm}

        \vspace{3pt}
        $x$-intercept: \blank{1.8cm} \quad $y$-intercept: \blank{1.8cm} \quad
        $f$ is increasing on \blank{2.6cm}

        \vspace{3pt}
        Absolute \blank{1.6cm} of \blank{1.0cm} at $x=\blank{1.0cm}$; \; no absolute
        \blank{1.6cm}.

  \item \textbf{Explain.} Why does $f$ have a restricted domain, while $\sqrt[3]{x+1}-2$ would not?
        And what happens to this graph as $x\to-\infty$?
        \writelines{3}

  \item \textbf{SOL-style multiple choice.} What is the domain of $h(x) = \sqrt{5-2x}$?
        \begin{enumerate}[label=\Alph*., leftmargin=1.9em, itemsep=1pt]
          \item $x \ge \dfrac{5}{2}$
          \item $x \le \dfrac{5}{2}$
          \item $x \ge -\dfrac{5}{2}$
          \item all real numbers
        \end{enumerate}
\end{enumerate}

\end{document}
